The AT command responses identified the radio firmware and confirmed its P2P work mode. The version query required the \texttt{=?} suffix, while the serial-number query returned no value. We recommend using the query syntax documented for the installed firmware and preserving the complete command-and-response output.

P2P activity was observed at both the transmitter and receiver. The longer receiver record contains 108 samples with varying RSSI and SNR. These measurements describe reception during the recorded interval, but do not determine maximum range or delivery ratio. For a range comparison, we recommend recording distance and obstacles alongside the receiver timestamps and preserving a transmitted-message count.

The MEGA--RAK setup transmitted an incrementing counter encoded as hexadecimal text. The video records matching counter payloads at the transmitter and receiver, supporting reception during the demonstration. We recommend retaining synchronized logs from both ends and identifying the exact sketch and radio configuration associated with each record, since the shared sketch and the code visible in the video use different spreading factors.
